\section{为什么先把体系切回 LaTeX}
一个稳定的 PDF 生产系统首先依赖版心、字体、行距、标题体系、分页算法与图表规则，而不是依赖大量“卡片组件”。LaTeX 的优势在于这些基础排印机制已经成熟，尤其适合中文长文、公式、表格和图形混排。

\BenchCallout{本轮基准要求}{正文应成为页面主体；中文行长与行距必须舒适；标题层级明确；公式与图表不破坏阅读节奏；任何视觉增强都不能替代排印基础。}

\section{同一内容的复杂度测试}
\subsection{正文与公式}
在一个机械系统中，动量守恒写作
\[
 m_1v_1+m_2v_2=m_1v'_1+m_2v'_2.
\]
如果碰撞后两物体粘在一起，则可进一步写成完全非弹性碰撞的共同速度表达式。这一段故意包含较长中文句子，用来观察行长、标点挤压、行距以及段落密度是否自然。

\subsection{表格}
\begin{center}
\begin{tabularx}{.94\linewidth}{@{}lXX@{}}
\toprule
项目 & 合格表现 & 常见失败 \\
\midrule
正文 & 连续阅读自然、行长稳定 & 像网页卡片，视觉跳跃 \\
公式 & 与正文基线协调 & 公式上下留白失衡 \\
图形 & 能与解释文字互相指代 & 图大于内容、抢夺注意力 \\
\bottomrule
\end{tabularx}
\end{center}

\subsection{图形}
\begin{center}
\begin{tikzpicture}[>=stealth,thick]
\draw[->](0,0)--(6.5,0) node[right]{时间};
\draw[->](0,0)--(0,3) node[above]{理解深度};
\draw[accent,very thick] plot[smooth] coordinates{(.4,.4)(1.5,.9)(2.8,1.5)(4.2,2.2)(5.8,2.6)};
\node[accent] at (4.3,1.55){结构化理解};
\end{tikzpicture}
\end{center}

\section{验收观察}
这一部分用于观察标题、正文、公式、表格、图形、页眉页脚与留白的整体节奏。模板通过的前提不是“看起来像设计稿”，而是连续阅读仍然舒适。

\BenchCallout{视觉检查清单}{中文正文不拥挤；标题层级明确但克制；图表与正文互相指代；页眉页脚不抢视觉焦点；长内容能够自然分页。}
